\mysection{Conclusion}\label{sec:conclusion}
In conlusion, we have performed linear feedback cooling to a magnetically levitated particle of submillimeter scale in a superconducting trap. Two different translational modes (the y- and x-mode) have simultaneously reached \SI{7.1}{mk} and \SI{6.6}{mK}. This is performed in a dry dilution refrigerator, with extensive vibration isolation. The mass of the levitated particle is \SI{0.38}{mg} and the feedback cooled phonon number is \SI{2.9e6}{} and \SI{2.0e6}{}, for mode 3 and 4 respectively.

Currently, we are limited by detection sensitivity and upconversion of resonances of the vibration isolation system below \SI{10}{Hz}. Improvements on the achieved mode temperatures in this experiment can be achieved with lower starting temperatures and/or lower force noise (i.e.  better vibration isolation), higher Q factors (which will also require an even better vibration isolation) and lower detection noise (which can be achieved by a larger coupling, i.e. placing the pick-up coil closer to the magnet). 

We can calculate the minimal achievable mode temperature and the corresponding phonon number of the current setup, due to force and detection noise:
\begin{comment}
\begin{equation}
    T_\text{min} = \sqrt{\frac{\omega_0}{2 \kB} \frac{k}{Q_0} S_{x_d}}
\end{equation}
\end{comment}
\begin{equation}
    T_\text{min} = \frac{\omega_0}{2 \kB} \sqrt{ S_F S_{x_\text{det}}}
\end{equation}

Reading from both Fig.~\ref{fig:cooling3} and Fig.~\ref{fig:cooling4} a position noise of $\SI{1}{pm/\sqrt{Hz}}$, and using the uncooled ($\text{gain} =0$) mode temperatures for the force noise, for mode 3 this results in $T_\text{min, 3} = \SI{0.65}{mK}$ and $N_\text{ph, min, 3} = \SI{2.7e5}{}$, and for mode 4 this results in $T_\text{min, 4} = \SI{1.9}{mK}$ and $N_\text{ph, min, 4} = \SI{5.7e5}{}$ as the minimal achievable values with the current setup. 
% mode 3 T: https://www.wolframalpha.com/input?i=%282*pi*50.59Hz%29%2F%282*boltzmann+constant%29+*+sqrt%284+boltzmann+constant*1926mK+*0.35mg*2*pi*50.59Hz%2F3.77e6*+%281pm%2Fsqrt%28Hz%29%29%5E2%29
%mode 3 phonons: https://www.wolframalpha.com/input?i=1%2F%282*hbar%29+*sqrt%284+boltzmann+constant*1926mK+*0.35mg*2*pi*50.59Hz%2F3.77e6*+%281pm%2Fsqrt%28Hz%29%29%5E2%29
%mode 4 T: https://www.wolframalpha.com/input?i=%282*pi*67.98Hz%29%2F%282*boltzmann+constant%29+*+sqrt%284+boltzmann+constant*9581mK+*0.35mg*2*pi*67.98Hz%2F5.45e6*+%281pm%2Fsqrt%28Hz%29%29%5E2%29
%mode 4 phonons: https://www.wolframalpha.com/input?i=1%2F%282*hbar%29+*sqrt%284+boltzmann+constant*9581mK+*0.35mg*2*pi*67.98Hz%2F5.45e6*+%281pm%2Fsqrt%28Hz%29%29%5E2%29

Keeping the current detection sensitivity, but improving the vibration isolation such that upconversion is no longer an issue and the force noise is determined by the thermodynamic temperature of the environment (\SI{20}{mK}), the phonon numbers would fall to \SI{2.7e4}{} for mode 3 and \SI{2.6e4}{} for mode 4. The current magnet, at a frequency of \SI{50}{Hz}, has a zero point motion of \SI{0.7}{fm}, this corresponds to a mode temperature of \SI{2.4}{nK}.
% mode 3: https://www.wolframalpha.com/input?i=1%2F%282*hbar%29+*sqrt%284+boltzmann+constant*20mK+*0.35mg*2*pi*50.59Hz%2F3.77e6*+%281pm%2Fsqrt%28Hz%29%29%5E2%29
% mode 4: https://www.wolframalpha.com/input?i=1%2F%282*hbar%29+*sqrt%284+boltzmann+constant*20mK+*0.35mg*2*pi*67.98Hz%2F5.45e6*+%281pm%2Fsqrt%28Hz%29%29%5E2%29
% x_zpm: https://www.wolframalpha.com/input?i=sqrt%28hbar+%2F+%282+*+0.35+mg+*+2+*+pi+*+50Hz%29%29
% T_zpm: https://www.wolframalpha.com/input?i=hbar+*+2+*pi+*+50+Hz+%2F+boltzmann+constant

In a more extreme case where the setup would be pushed towards its limits, with a thermodynamic temperature of \SI{2}{mK}~\cite{Loek, DeWitVibIso}, a 10 times higher Q factor, a detection sensitivity of $\SI{0.01}{pm/\sqrt{Hz}}$, we would be able to cool mode 3 to \SI{27}{} phonons and mode 4 to \SI{26}{} phonons.
% mode 3: https://www.wolframalpha.com/input?i=1%2F%282*hbar%29+*sqrt%284+boltzmann+constant*2mK+*0.35mg*2*pi*50.59Hz%2F3.77e7*+%280.01pm%2Fsqrt%28Hz%29%29%5E2%29
% mode 4: https://www.wolframalpha.com/input?i=1%2F%282*hbar%29+*sqrt%284+boltzmann+constant*2mK+*0.35mg*2*pi*67.98Hz%2F5.45e7*+%280.01pm%2Fsqrt%28Hz%29%29%5E2%29

%When increasing the measurement coupling
Further improving the experiment and lowering the phonon amount towards the ground state, the back action noise from the SQUID will dominate over the thermal noise of the experiment, at some measurement coupling~\cite{ULDM}. This calls upon improving the SQUID noise to the quantum limit, since the SQUID we use has an energy resolution of approximately 10.000 times $\hbar$~\cite{ULDM, Ahrens2025, Vinante2021}.

Combining all these improvements, cooling and detection of the levitated particle in the quantum ground state should become possible.

% Furthermore, if the current levitated system is seen as a quantum system, the coherence time can be calculated. During the experiment the bath temperature (temperature of the trap and its surroundings) is typically \SI{20}{mK}. At a resonance frequency of \SI{50}{Hz} this corresponds to a bath phonon occupation of \SI{4.2e8}{}. At this frequency, the ringdown time is \SI{2.4e4}{s}, and the coherence time ($\tau_\text{coh}$) of the system would be $\tau_\text{coh} = \tau_\text{ringdown}/N_\text{ph} = \SI{58}{us}$. \DU{hier nog die andere twee getallen bijzetten}

%\DU{@Tjerk, je had het over 1K is equivalent met 20GHz. Ik kom daar op uit als je neemt $\kB T = hf$. Ik zou zelf denken dat je $1/2 \kB T= hf$ moet nemen.} Tjerk: je hebt twee vrijheidsgraden positie 1/2 kx^2 en snelheid 1/2mv^2. Die hebben allebei 1/2 kBT, dacht ik. N_phonon = k_BT/hf, dacht ik.


\begin{comment}
\vspace{2cm}
To do:
\begin{itemize}
%    \item nadruk leggen op dat het phonon nummer wel hoog is maar dat dit systeem al gebruikt is voor zwaartekracht
%    \item ergens quantum limited squid en betere vib iso (upconversion) noemen. Dan zien we 1fm meten wel zitten
%\item  upconversion noemen in results. Zie hieronder?
%\item DU twijfelt uit zijn hoofd hoe we nu de koeling uitleggen, volgens mij noemen we de welch method e.d. nog nergens
\item
\end{itemize}


    x_\text{zpm} = \sqrt{\frac{\hbar}{2 m \omega_0}}

\DU{met een 10 keer hogere koppeling, dus 100 keer hogere m/V, en een quantum limited squid (10.000). Dan nog start T 100x lager, dan ben je er al. Maar dan moet je de upconversion bestreiden. Dus de huidige micrometers moeten nanometers 0-15Hz. Dus betere vib iso. Liefst ook de vib iso modes in een kleiner bereik dan de huidige 0-15Hz.\\


\end{comment}
